\begin{recipe}
[%
    preparationtime = {\unit[45]{min}},
    portion = {\portion{4}},
    source = {\protect\recipesource{https://www.instagram.com/p/DTBH5jllSbC/}{FIT FOR FUN\: Veganes Butter Chicken}},
    calory = {400kcal (700kcal) | 22g Protein (30g) | 31g KH (86g) | 18g Fett (23g)},
    price = {1,80€ (2,40€)},
    created = {28.07.2026},
    rating = {3}
]
{Veganes Butter Chicken mit Sojamedaillons}

\graph{% pictures
    big=pic/hauptgericht/21_veganes-butter-chicken
}

\introduction{%
Kräftig gewürzte Sojamedaillons in einer cremigen Tomatensauce mit Garam Masala, Ingwer und einem Hauch Zimt. Die vegane Variante bleibt nah am Charakter eines Butter Chicken, bringt aber ordentlich pflanzliches Protein mit.
}

\ingredients{
    \textbf{Sojamedaillons} & \\
    2 Handvoll & Sojamedaillons\index{Tofu \& Fleischalternativen!Soja}\\
    2 TL & Garam Masala\\
    2 TL & Chilipulver\\
    2 TL & Paprikapulver, edelsüß\\
    2 TL & Kreuzkümmel\\
    & Salz, Pfeffer\\
    & Gemüsebrühe zum Einweichen\\
    & Kokosöl zum Braten\\
    \textbf{Sauce} & \\
    2 & Zwiebeln\\
    1 Stück & Ingwer\index{Gemüse!Ingwer}\\
    1 Zehe & Knoblauch\index{Gemüse!Knoblauch}\\
    2 EL & Vegane Butter\index{Milchprodukte \& Alternativen!Vegane Butter}\\
    2 TL & Paprikapulver, edelsüß\\
    2 TL & Garam Masala\\
    2 TL & Kreuzkümmel\\
    2 Dosen & Gestückelte Tomaten\index{Gemüse!Tomate}\\
    200 ml & Pflanzliche Sahne\index{Milchprodukte \& Alternativen!Pflanzliche Sahne}\\
    3 EL & Sojajoghurt\index{Milchprodukte \& Alternativen!Sojajoghurt}\\
    1 Prise & Zimt\\
    & Salz, Pfeffer
}

\preparation{%
    \step%
    Gemüsebrühe aufkochen, Sojamedaillons hineingeben und bei geschlossenem Deckel etwa 15 Minuten ziehen lassen. Anschließend abgießen und die Medaillons vorsichtig ausdrücken.

    \step%
    Medaillons von beiden Seiten mit Garam Masala, Chili, Paprika, Kreuzkümmel, Salz und Pfeffer würzen.

    \step%
    Etwas Kokosöl in einer großen Pfanne erhitzen. Die Medaillons darin von beiden Seiten jeweils 4--5 Minuten kräftig anbraten, bis sie gut gebräunt sind. Danach aus der Pfanne nehmen.

    \step%
    Zwiebeln, Ingwer und Knoblauch schälen und fein hacken. Erneut etwas Kokosöl in die Pfanne geben und alles zusammen mit Paprikapulver, Garam Masala und Kreuzkümmel bei mittlerer Hitze anschwitzen, bis die Zwiebeln glasig sind und die Gewürze duften.

    \step%
    Gestückelte Tomaten, pflanzliche Sahne, Sojajoghurt und Zimt einrühren. Die Sauce etwa 10 Minuten bei mittlerer Hitze köcheln lassen.

    \step%
    Sauce fein pürieren. Sojamedaillons zurück in die Pfanne geben und die vegane Butter einrühren. Noch einmal kurz erwärmen und mit Salz, Pfeffer und bei Bedarf etwas Garam Masala abschmecken.
}

\hint{
    Die eingeweichten Medaillons nur vorsichtig ausdrücken. Sind sie zu nass, bräunen sie schlecht; drückt man sie komplett trocken, werden sie später schnell zäh.
}

\end{recipe}
